\bookStart{Lay of Hymer}[Hymis kviða]
\setBookCode{Hymiskvida}

\begin{flushright}%
\textbf{Dating} \parencite{Sapp2022}: C10th (0.694)

\textbf{Meter:} \Fornyrdislag%
\end{flushright}%

\section{Introduction}

The \textbf{Lay of Hymer} (\textlink{Hymiskvida}) comedically narrates how the gods Thunder and Tew go on an adventure into Ettinland in order to fetch a great cauldron, and how Thunder pulls up the Middenyardswyrm with his fishing pole.

\subsection{Preservation}

\Hymiskvida\ attested in both \Regius\ and \AM.  The two mss. agree very well with each other; they share the same stanzas in the same order.  The most substantial difference is the title; \AM\ has \emph{Hymis kviða} ‘the lay of Hymer’ while \Regius\ instead has \emph{Þórr dró Miðgarðs-orm} ‘Thunder pulled the Middenyardswyrm’.

\subsection{The story}

\textlink{Hymiskvida} belongs to a genre of comedic stories about Thunder’s adventures in Ettinland.  These stories all see Thunder set out from Osyard followed by one or more companions.  Together the party undergoes various fantastical tests and hardships on their way through the Ettinhomes, and return home after Thunder has slaughtered yet another Ettin clan.  The popularity of this genre is seen by the large number of stories belonging to it.  In the Poetic Edda is also represented by \textlink{Thrymskvida} and to some extent \textlink{Harbardsljod}.  Other examples are Thunder’s journey to Outyards-Lock in \Gylfaginning~44--47, his fight with Thedse in \Skaldskaparmal\ 2--3 and \Haustlong\ 2--13, his fight with Rungner in \Skaldskaparmal\ 24 and \Haustlong\ 14--20, and his journey to Garfrith in \Skaldskaparmal\ 26 (in which two sts.~are cited; see \textlink{EddicFragments}[F6][VI]).

\subsubsection{The otherness of the Ettins}

In these stories the Ettins are very much an \emph{other} to the Gods, and this is something which \textlink{Hymiskvida} strongly emphasizes:

\begin{itemize}
  \item The Ettins live in the far east (st.~5) in an inhospitable, frozen climate (st.~10) of mountains (sts.~2, 17) and lavafields (sts.~36, 38);
  \item they are physically deviant: misshapen (st.~10), grey-haired (st.~16), many-headed (sts.~8, 35), having heads harder than stone (sts.~30--31);
  \item they are likened to apes (st.~20), whales (st.~36) and Danes (st.~17, see note!);
  \item they are stingy and inhospitable (sts.~9, 16);
  \item they are snide and cowardly (sts.~19--20, 25--26, 28--32).
\end{itemize}

These traits of the Ettins stand in direct opposition to the Old Germanic social norms as represented and upheld by the Gods, who live in a lush green land and are young, beautiful, generous, and brave.  The one exception in the poem is Tew’s mother in st.~8, who is blonde, beautiful and hospitable.  The mother of a god should also be godlike, and this likely reflects the broader motif of a beautiful ettin-woman who is taken by a god, for which cf.~the st.~8 n.~and \textlink{Skirnismal} above.

As natural inferiors and a threat to the social order the Ettins must be subjugated by the Gods, and the agent of this is Thunder.  Throughout the poem he recurringly humiliaties the ettins Eagre and Hymer through completing their various challenges and performing tests of strength.  These tests follow a similar outline: Thunder is given a dangerous or near-impossible challenge but accomplishes it anyway with surprising ease through a combination of brawn and brain to the humiliation of the challenger.  The challenges consist of finding an enormously large kettle (st.~3, explicitly called Eagre’s “revenge”), wrestling one of Hymer’s oxen for bait (sts.~17--18), carrying home Hymer’s whales and boat (st.~26), breaking Hymer’s finest chalice (st.~28), and perhaps also taking away the cauldron (st.~33)---though that may just be Hymer wishing to finally be rid of the pestering gods.

In the end Thunder delivers justice by slaughtering Hymer and his whole clan of many-headed Ettins.

\subsubsection{The fishing expedition}

At the center of the poem stands Thunder’s famous fishing trip, where the god gets the Middenyardswyrm on his hook but ultimately fails to catch it.  In this section of \Hymiskvida, especially in sts.~22--24, we find a noticably more reverent tone than elsewhere in the poem.

Thunder’s abortive encounter with the Wyrm was clearly very popular in the Wiking Age and is, apart from \Hymiskvida, also dealt with in five fragmentary Scaldic poems from the 9th or 10th centuries and extensively in \Gylfaginning~48:3--31.  The story is also represented pictorally on numerous objects from the Wiking Age.

The Scaldic fragments are all found in citations scattered throughout \Skaldskaparmal, all edited in \Skp\ 3.  They are, by descending number of long-lines: Bragi \emph{Þórr}, \Husdrapa\ 3--6, EVald \emph{Þórr}, and Ggnæv \emph{Þórr}, Ǫlv \emph{Þórr}.

In their present state none of these fragments are complete narratives, but mainly focus on the scene of Thunder in the boat facing off against the hooked Wyrm pressed against the gunwale.
From what little we can glean from them they also disagree on the course of events.  In Bragi \emph{Þórr} 6 the encounter ends when the cowardly Hymer cuts the fishing line and the Wyrm sinks back unscathed into the sea in \Husdrapa\ 6 and probably also Ggnæv \emph{Þórr} 1 Thunder strikes the head off the Wyrm with his hammer Millner.

Other than \textlink{Hymiskvida} the only complete retelling of the myth is found in \Gylfaginning~48:3--31, which may be summarized as follows:

\begin{quote}{\small Thunder goes out alone into Middenyard in the shape of a young man (\emph{ungr dręngr}) without his goats and chariot.  In the evening he comes to the ettin Hymer and asks to stay the night.  At dawn Hymer plans to go fishing and Thunder asks to join him.  Hymer says that Thunder will not be of much help since he is “small and but a young man” (\emph{lítill ok ung-męnni ęitt}), and warns him that he usually takes long and arduous trips.  Thunder, angered, says that he will row very far, and then asks Hymer what bait they will use.  Hymer tells him to find it himself, and so he turns to his flock of oxen; he tears off the head from the greatest ox, one called Heavenrid (\emph{Himin-hrjóðr}).
The two go out to sea, and Thunder rows far past Hymer’s usual fishing waters.  Hymer, unhappy, warns him that if they row any further out they will be in danger of the Middenyardswyrm, but Thunder keeps on.  After some time he puts down the oars, readies his fishing line, hooks the ox-head and lowers it.  The Wyrm soon bites, and struggles so hard that Thunder is pressed against the gunwale.  In rage he brings himself into his Os-might (\emph{ǫ́s-męgin}) and pulls back with such force that his feet go through the bottom of the ship and press into the seabed.  The Wyrm's head goes up against the gunwale.  The two enemies ferociously stare at each other, Thunder “sharpening his eyes” and the Wyrm spitting venom.  Hymer is frightened, reaches for his bait-cutting knife, and cuts the line---the Wyrm then sinks back into the sea.  Thunder throws his hammer after it, “and men say that he struck off the monster’s head, but I think it true to tell thee that the Middenyardswyrm still lives and is lying in the outer sea.”  Thunder gives Hymer a punch to the ear so that he flies headfirst overboard; the god then wades back to land.}\end{quote}

Snorre’s account is clearly based on multiple sources, certainly including the Scaldic fragments he himself cites in \Skaldskaparmal.  It is hard to say whether Snorre had access to \Hymiskvida; the closest verbal agreement is when it is said in \Gylfaginning~48:22 that \emph{Mið·garðs-ormr gein yfir uxa-hǫfuð’it, en ǫngull’inn vá í góm’inn orm’inum} ‘The Middenyardswyrm snapped at the ox-head and the hook went into the roof of the wyrm’s mouth’, which has some resemblance to st.~22.

Some details in \Gylfaginning~48 not found in \Hymiskvida\ must derive from now-lost texts available to Snorre; most famously the detail of Thunder’s feet going through the boat is also found on the Swedish Altuna stone and the Danish Hørdum stone (but see note to st.~34/2 below).  The name of the ox Heavenrid is attested in \inx[C]{thule}[thules] listing names of oxen.

In addition to the literary sources there are also numerous pictorial depictions of the myth from the Wiking Age.  These are the Swedish runestones from Altuna (U 1611, figure~\ref{fig:altuna}, p.~\pageref{fig:altuna}) and Linga (Sö 352), several Jutlandic picture stones from Hørdum, a Cumbrian picture stone from Gosforth, and the Gotlandic picture stone GP 21 from Ardre church.  The images depict the same scene as the Scaldic fragments: Thunder stands in the boat above the hooked Wyrm, often depicted as a fish; next to him is one companion.  Some of them have additional details like the use of the ox-head for bait (U 1611, Sö 352), Thunder’s foot going through the boat (U 1611, Hørdum), or Thunder raising his hammer (U 1611).

More broadly, Thunder’s fishing reflects the archetypal fight (the \emph{Chaoskampf}) between the Storm-God and the Dragon, as found in a great many mythologies.  Important examples of this include Vedic \textoi{Índra} and \textoi{Vr̥trá} (\Rigveda\ 1.32 et c.), Babylonian \emph{Marduk} and \emph{Tiamat} (\emph{Enūma Eliš} TODO), Greek \textgreek{Ζεύς} (\emph{Zeús}) and \textgreek{Τῡφῶν} (\emph{Tūphôn}), Hebrew YHWH and Leviathan (TODO: references).  With these analogies in mind it seems that the versions where Thunder slays the Wyrm reflect an older layer of Germanic mythology, before the lethal fight between Thunder and the Wyrm had been transposed to the End Times (see \textlink{Voluspa}[53]).

\subsubsection{\Hymiskvida\ as a composite}

The narrative of \textlink{Hymiskvida} can be subdivided into the following episodes:

\begin{enumerate}
  \item Sts.~1--6: The Gods wish to drink ale and Thunder goes to Eagre to force him to brew it; Eagre in turn asks for a cauldron big enough to brew enough ale for all the Gods.
  \item 7--16: Thunder and Tew go to visit Tew’s father, the stingy ettin Hymer, who owns such a cauldron; horrified at Thunder’s great appetite during the evening he tells them that they must go fishing for food.
  \item 17--19: Thunder says that he will do it, if he is given bait; Hymer challenges him to kill one of his oxen; Thunder tears off the head from one of them.
  \item 20--25: The three go fishing; Hymer pulls up some whales; with the ox-head as bait Thunder manages to hook the Middenyardswyrm itself; he loses it.
  \item 26--27: Hymer challenges Thunder to carry the boat and whales back to his farm; he does.
  \item 28--32: Hymer challenges Thunder to break a supposedly indestructible chalice; he succeeds by smashing it against the ettin’s forehead.
  \item 33--36: Thunder and Tew depart with the cauldron; they find themselves followed by Hymer and his ettins; Thunder kills them all.
  \item 37--38: One of Thunder’s goats goes halt.
  \item 39: Thunder returns to the Gods with Hymer’s cauldron; they host a banquet.
\end{enumerate}

The fishing expedition found in the Scaldic fragments and \Gylfaginning~48 is represented by episodes 3--4.
\textlink{Hymiskvida} is the only source to place the fishing expedition within the broader context of Thunder and Tew obtaining a cauldron from Hymer for the sake of brewing ale (episodes 1--2, 9).

It seems inescapable, both on the basis of the other Norse sources and broader comparative mythology, that both the frame narrative and the additional episodes (e.g.~number 8) originally had nothing to do with Thunder’s encounter with the Wyrm.  They have in any case been woven together into a single narrative, probably by the poet himself, for the sake of a more entertaining story for an audience which was likely expecting certain common “Thunder’s journey” tropes.

Although the poem is skilfully written this weaving-together has not been entirely successful and has resulted in a few few loose threads.  The halt goat of sts.~37--38 (episode 8) finds a parallel in \Gylfaginning~44:17--21 where it serves as the origin story of Thunder’s two child servants Thelve and Wrash who accompany him on the way to Outyards-Lock.  These two children play an important part in the narrative of \Gylfaginning~44--47, but in \textlink{Hymiskvida} are an entirely superfluous detail---something the poet himself anticipates in his address to the audience.  It is also strange that Lock should appear so late in the poem when he is not mentioned before or after; he too plays an important role in the Outyards-Lock story of \Gylfaginning~44--47.

A major superfluous element is the god Tew, who plays no role at all in the fishing expedition.  He is last alluded to in st.~16 where Hymer speaks of “[us] three”, only to reappear in st.~33 where he fails to lift the cauldron.  It is also rather strange that Tew shows no reaction to the murder of Hymer, his father, taking place right before his eyes (sts.~35--36), but for the paternity cf.~\textlink{Hymiskvida}[5]/3 n.

The simplest explanation for this is however that Tew originally had nothing to do with the encounter between Thunder and the Wyrm, and within the confines of \textlink{Hymiskvida} mainly serves to segue from the narrative of the cauldron-fetching to the fishing expedition.  Strong support for this is found in the fact that Thunder in the other variants of the fishing-story only has one companion---Hymer.  These variants include the pictoral depictions, which only show two figures seated in the boat.

\subsection{Style}

When speaking of a composite poem, one must distinguish between a text where several originally separate works have been put together mostly unchanged, and a text composed by a single author drawing from multiple narrative sources.  A likely example of the former is \textlink{Havamal}, but \textlink{Hymiskvida} undoubtedly belongs to the latter category.

\Hymiskvida\ has a distinct style and rhythm throughout which stands alone in the Poetic Edda.  Indeed, the sharpest stylistic contrast is with the poem most similar content-wise, \textlink{Thrymskvida}.  Where \textlink{Thrymskvida} is written in a rustic style with fairly loose \Fornyrdislag\ meter, repeats lines and even whole stanzas, \textlink{Hymiskvida} instead uses an unusually strict meter, avoids repetition, and revels in kennings, difficult grammatical constructions, and highly unnatural word order (see especially sts.~16/3--4, 20, and 39/3--4).  These are all traits one associates more closely with Scaldic poetry in complex measures like \Drottkvett\ than with Eddic poetry in \Fornyrdislag, and it seems certain that the anonymous poet of \textlink{Hymiskvida} must have had some training in the Scaldic art and familiarity with compositions in that genre.  Two kennings (17/4a \emph{brjótr berg-Dana}, 22/4 \emph{umb-gjǫrð allra landa}) are even shared identically with Scaldic poems in \Drottkvett.

\subsection{Formal characteristics}

\Hymiskvida\ consists of 152 long-lines divided into 39 stanzas.  The typical length of a stanza is 4 lines (32/39 sts.: 1--5, 7--10, 12--23, 28--35, 37--39).  Deviant stanzas with 2 lines (1 st.: 6), 3 lines (4 sts.: 24--26, 36), and 5 lines (2 sts.: 11, 27) are also represented.  The 4-line stanzas generally have a semantic division between ll.~2 and 3 and thus consist of two helmings.

Of the deviant stanzas, st.~6 appears to represent an isolated helming and is thus likely missing two lines.
Two out of four 3-line sts.~(24 and 25) are clearly corrupt owing to syntactical issues; on the other hand the two others (26 and 36) show no obvious corruption.
The two 5-line stanzas do not seem to be missing any lines; in both of them one may identify an additional fifth line (11/1, 27/5) which could be removed without hindering the form of the other 4, although that is not to say that these lines are necessarily later inserts.  To conclude, the original \Hymiskvida\ most likely only tolerated 4-line sts., with occasional exceptions for 5-line stanzas with a marked fifth line.

The meter of \textlink{Hymiskvida} is \Fornyrdislag, but of a stricter variant than that of any other poem in the Poetic Edda; this is especially true when it comes to the count and weight of syllables.  Three-syllable lines are rare, and the ones that do exist can all be corrected by restoring older hiatus forms (see note to 23/1b below).  The poet also has a notable tendency to avoid lines of type C1 (xP\#Px) where both stresses are heavy, something shared with Scaldic poems in \Kviduhattr\ and with Eyel’s Headransom.  This probably explains his unique tendency to place the two-syllable preposition \emph{fyrir} ‘before, in front, (up) ahead’ at the end of the b-verse (never the a-verse), which he does 6 times---more frequently than in any other \Fornyrdislag\ poem of the Poetic Edda.

\begin{figure}[h]
\centering
\includegraphics[width=\textwidth]{Altuna}
\caption{Detail of Thunder on the Altuna stone.  Wiking Age, ca. 1000 \ce.  © Bengt A. Lundberg, RAÄ. \href{https://creativecommons.org/licenses/by/4.0/deed.en}{CC BY 4.0}.  \url{https://pub.raa.se/dokumentation/a4396e83-dac1-4f69-9cab-ff877da64e8d}}
\label{fig:altuna}
\end{figure}

\newpage

\section{Text}

\subsection{The Lay of Hymer}

\bvg\bva\mssnote{\Regius~13v/26, \AM~5v/25}%
Ár \edtrans{\alst{v}al-tívar}{slain-Tews}{\Bfootnote{The Gods, a rare epithet only otherwise found in \textlink{Voluspa}[50]/2, 59/4 (n.).}} \hld\ \alst{v}ęiðar nǫ́mu &
ok \edtrans{\alst{s}umbl-\alst{s}amir}{assembled at the simble}{\Bfootnote{A simble being a drinking-feast or banquet.  For the simble of the Gods see \textlink{Grimnismal}[37]/6, 46/4--5, \textlink{Lokasenna} (\emph{passim}).}} \hld\ \edtrans{áðr \alst{s}aðir yrði}{before they might eat}{\Bfootnote{Lit.~‘might become sated’.}}, &
\edtrans{\alst{h}ristu tęina \hld\ ok ȧ \alst{h}laut sǫ́u}{they shook the twigs and looked upon the leat}{\Bfootnote{The Gods performed an augury, the details of which are not clear from this line alone.  The term \emph{hlaut} ‘leat’ is explained in \HakonarSaga\ and \EyrbyggjaSaga\ as the sacrificial blood of the slaughtered beasts, which was sprinkled by means of \emph{hlaut-tęinar} ‘leat-twigs’.  These two words are also found together in an OE glossary, where \emph{tân-hlyta} ‘twig-allotter’ glosses Latin \emph{sortilegus}.
Etymologically ON \emph{hlaut} is identical to OS \emph{hlôt} ‘lot’ (\textlink{Heliand}[66][5547b]), Gothic \emph{hlauts} ‘lot, inheritance’ \parencite[58]{Streitberg1910}.
If we trust these sources the simplest explanation is that the Gods sprinkled the animal blood and interpreted the patterns.  In any case they found it most auspicious to feast at Eagre’s.}}, &
fundu at \edtrans{\alst{Ę́}gis}{Eagre’s}{\Bfootnote{\emph{Ę́gir} ‘Eagre’ is the spirit of the Sea, an ettin.}} \hld\ \alst{ø}r·kost hvera.\eva

\bvb%
{\huge O}\textsc{f yore the slain-Tews} \name{Gods} had caught game, \\
and, assembled at the \inx[C]{simble} before they might eat, \\
they shook the twigs and looked upon the \inx[C]{leat}; \\
they found at Eagre’s a great choice of cauldrons.\evb\evg


\bvg\bva\mssnote{\Regius~13v/28, \AM~5v/27}%
\Ballnote{Having seen that Eagre has a great store of cauldrons, Thunder visits him and orders him to brew ale for the thirsty Eese.}%
Sat \alst{b}erg-\alst{b}úi \hld\ \alst{b}arn-tęitr fyrir, &
\edtrans{\alst{m}jǫk g·líkr \alst{m}ęgi \hld\ \alst{M}iskur·blinda}{much like the son of Misherblind}{\Bfootnote{An unexplained reference.  \emph{Miskur·blindi} ‘Misherblind’ might be another name for Firneet (\emph{Fornjótr}), the father of Eagre (Sea), Lay (\emph{Logi} ‘Fire’), and Wind (\emph{vindr}; \Skaldskaparmal\ 35, 36, \Ynglingatal).  The line would then be a tautology: “he looked just like the son his father”.}}, &
lęit ï \alst{au}gu \hld\ \alst{Y}ggs barn ï þrá: &%NOTE: þrá is correct; it is not the same word as þrǫ́
\llap{„}þú skalt \alst{ǫ̇}sum \hld\ \alst{o}pt sumbl \edtext{gøra}{\lemma{gøra ‘make’}\Afootnote{\emph{gefa} ‘give’ \AM}}!“\eva

\bvb The crag-dweller \ken*{\textsc{ettin} = Eagre} sat merry like a child ahead, \\
much like the son of Misherblind. \\
Into his eyes looked Ug’s \name{Weden’s} child \ken*{= Thunder} defiantly: \\
“Thou shalt for the Eese oft make simbles!”\evb\evg


\bvg\bva\mssnote{\Regius~13v/31, \AM~5v/29}%
\Ballnote{Eagre is upset at the demand but too cowardly to openly challenge Thunder.  He instead gets back at him by demanding a single cauldron big enough to supply ale for all the Eese.}%
\alst{Ǫ}nn fekk \alst{jǫ}tni \hld\ \alst{o}rð-bę́ginn halr, &
\alst{h}ugði at \alst{h}efndum \hld\ \alst{h}ann nę́st við goð, &
bað \alst{S}ifjar ver \hld\ \alst{s}ér fǿra hver, &
\llap{„}þann’s ek \alst{ǫ}llum \alst{ǫ}l \hld\ \alst{y}ðr of hęita.“\eva

\bvb Great toil for the ettin the word-peevish man \ken*{= Thunder} caused; \\
he thought of revenge, soon, against the god. \\
He bade Sib’s husband \ken*{= Thunder} bring him a cauldron, \\
“that one with which I for you all might warm ale.”\evb\evg


\bvg\bva\mssnote{\Regius~14r/1, \AM~5v/30}%
Né þat \alst{m}ǫ́ttu \hld\ \alst{m}ę́rir tívar &
ok \alst{g}inn-ręgin \hld\ of \alst{g}eta hvęr⸗gi, &
und’s af \alst{t}ryggðum \hld\ \alst{T}ýr \edtrans{Hlór·riða}{Loride}{\Bfootnote{‘The Loud/Bellowing Rider’, a name for Thunder.}} &
\alst{ǫ̇}st-ráð mikit \hld\ \alst{ęi}num sagði:\eva

\bvb Yet of one could not the famous \inx[P]{Tews} \\
and the \inx[P]{Yin-Reins} anywhere get hold--- \\
until, out of loyalty, Tew to Loride \\
a great loving counsel in private told:\evb\evg


\bvg\bva\mssnote{\Regius~14r/3, \AM~6r/2}%
\llap{„}Býr \edtext{fyr \alst{au}stan \hld\ \alst{É}li-vága}{\lemma{fyr austan Éli-vága ‘To the east of the Ilewaves’}\Bfootnote{According to \textlink{Vafthrudnismal}[31] the Ilewaves were the poisonous wild rivers from which the ettins emerged, and so it makes sense that they would be found in the east, where the ettins dwell.  That Hymer should dwell even to the east of the Ilewaves then illustrates his unusually ettin-like nature.

For the notion we may compare the Greek \textgreek{Ὠκεανός} ‘Ocean’, the world-encircling river, beyond which the world is supernatural and uncanny (\Odyssey\ 11.13--22, 155--160).  In \Theogony\ 790--1, \textgreek{Ὠκεανός} consists of nine rivers which flow out into the world-sea that surrounds the Earth (cf.~\textcite[201]{West1966}), which has a Norse parallel in the Outer Sea (\emph{Út-haf}), the sea surrounding the Earth in which the Middenyardswyrm dwells (cf.~l.~2 n., 22/4 n.~below).  Like the Middenyardswyrm which bites its tail, so is \textgreek{Ὠκεανός} given the epithet \textgreek{ἀψορρόος} ‘back-flowing, flowing into itself’ (\Iliad\ 18.399, \Odyssey\ 20.65, \Theogony\ 776).}} &
\edtrans{\alst{h}und-víss}{hundred-wise}{\Bfootnote{Alternatively “hound-wise”; the prefix simply means “very”, but the word is used exclusively for ettins, typically in the collocation \emph{hund-víss jǫtunn} (so \textlink{HelgakvidaHjorvardssonar}[25], \Skaldskaparmal\ 24, 26, \GautreksSaga\ 3, et c.).}} \alst{H}ymir \hld\ at \edtrans{\alst{h}imins ęnda}{at heaven’s end}{\Bfootnote{The heaven is thought of as a solid dome of stone, the earth as a flat disc (\textlink{Voluspa}[3] n., \textlink{Harbardsljod}[15] n.)  Hymer thus lives at the very edge of the disc, next to the Outer Sea which encircles all lands and in which the Wyrm lies.
For another ettin living at the edge of heaven cf.~\textlink{Vafthrudnismal}[37].}}; &
á \edtrans{\alst{m}ïnn faðir}{my father}{\Bfootnote{A notable detail since Tew is elsewhere called the son of \inx[P]{Weden} (\Skaldskaparmal\ 16) and it would be rather strange for the son of a male ettin to be afforded a place among the Eese due to the Norse patriarchal concern for male lineage.  A possible explanation for this discrepancy is if Weden slept with Tew’s mother, the beautiful woman of st.~8, which would be an instance of the motif of a god mating with a beautiful ettin-woman as seen e.g.~in \textlink{Skirnismal}.}} \hld\ \alst{m}óðugr kętil, &
\edtext{\alst{r}úm-brugðinn}{\Afootnote{\emph{†rumbrygðan†} \AM}} hver \hld\ \alst{r}astar djúpan.“\eva

\bvb “To the east of the \inx[P]{Ilewaves} dwells \\
the hundred-wise Hymer, at heaven’s end. \\
My father \ken*{= Hymer} owns, fierce, a kettle: \\
a roomy cauldron one \inx[C]{rest} deep.”\evb\evg


\bvg\bva\mssnote{\Regius~14r/4, \AM~6r/4}%
\Ballnote{Like elsewhere in this poem the speakers are not indicated, but it is most sensible that Thunder asks and Tew answers.}%
\speakernote{[Þȯrr kvað:]}%
\llap{„}Vęitst, ef \alst{þ}iggjum \hld\ \alst{þ}ann lǫg-velli?“ &
\speakernote{[Týr kvað:]}%
\llap{„}Ef, \alst{v}inr, \alst{v}élar \hld\ \alst{v}it gørum til!“\eva

\bvb\speakernoteb{[Thunder quoth:]}%
“Knowest thou if we’ll get that liquid-boiler \ken{cauldron}?” --- \\
\speakernoteb{[Tew quoth:]}%
“If, friend, we two make use of wiles!”\evb\evg


\bvg\bva\mssnote{\Regius~14r/5, \AM~6r/4}%
Fóru \alst{d}rjúgum \hld\ \edtrans{\alst{d}ag þann fram\emph{an}}{from the beginning of the day}{\Afootnote{emend.~after \textcite{FinnurEdda}; \emph{dag þann fram} ‘on that day forth’ \Regius; \emph{dag fráliga} ‘swiftly at day’ \AM}} &
\alst{Ǫ̇}s-garði frȧ \hld\ und’s til \edtrans{\alst{Ę}gils}{Eyel}{\Afootnote{so \Regius; \emph{Ę́gis} ‘Eagre’ \AM.}\Bfootnote{The reading of \AM\ is probably from confusion with the ettin Eagre, who is mentioned earlier in the poem.  Eyel, who takes Thunder’s goats in possession, is not otherwise known.  He may perhaps be identified with the farmer in \Gylfaginning~44:10--21, for which see Note to st.~37 below.}} kvǫ́mu; &
\alst{h}irði \alst{h}afra \hld\ \alst{h}orn-gǫfgasta; &
\alst{h}urfu at \alst{h}ǫllu \hld\ es \alst{H}ymir átti.\eva

\bvb They journeyed far from the beginning of the day \\
away from Osyard until they came to Eyel; \\
he took care the he-goats noblest of horns--- \\
they turned to the hall which Hymer owned.\evb\evg


\bvg\bva\mssnote{\Regius~14r/7, \AM~6r/6}%
\alst{M}ǫgr fann ǫmmu, \hld\ \alst{m}jǫk lęiða sér, &
\edtrans{\alst{h}afði \alst{h}ǫfða \hld\ \alst{h}undruð níu}{of heads she had nine hundred}{\Bfootnote{Malformed bodies, especially with a deviant number of body parts, are typical of ettins.  Other examples include a three-headed thurse in \textlink{Skirnismal}[31], the nine-headed ettin Thriwold (Bragi Frag 3 in \Skp\ 3), and the eight-armed Starked Eeldreng.  Cf.~Introduction and st.~35 below.}}, &
en \edtrans{\alst{ǫ}nnur}{another woman}{\Bfootnote{The use of the word “son” in the following line reveals this as Tew’s mother.  The poet stresses her beauty of dress and countenance, in contrast to the grandmother.}} gekk \hld\ \alst{a}l-gollin framm &
\alst{b}ru̇n-hvít \alst{b}era \hld\ \alst{b}jór-vęig syni:\eva

\bvb The lad \ken*{= Tew} found his grandmother very loathsome; \\
of heads she had nine hundred. \\
But another woman, all-golden, stepped forth, \\
white-browed, bringing a beer-draught for her son \ken*{= Tew}:\evb\evg


\bvg\bva\speakernote{[Týs móðir:]}\mssnote{\Regius~14r/9, \AM~6r/8}%
\llap{„}\alst{Á}tt-niðr \alst{jǫ}tna \hld\ \alst{e}k vilja’k ykkr &
\edtext{\alst{h}ug-fulla}{\Afootnote{\emph{hug-fulla hug-fulla} \Regius}} \edtrans{tvȧ}{two}{\Bfootnote{Here an instance of contraction < \emph{*tvȧa} also found in st.~15/4 below, but cf.~21/2b which has the disyllabic form.  It appears the poet was willing to use either form to accomodate the meter.}} \hld\ und \alst{h}vera sętja; &
es \alst{m}ïnn \edtrans{frí\emph{i}}{lover}{\Afootnote{so \Regius; \emph{faðir} ‘father’ \AM}} \hld\ \alst{m}ǫrgu sinni &
\edtext{\alst{g}løggr við \alst{g}ęsti \hld\ \alst{g}ǫrr ills hugar}{\lemma{gløggr \dots\ hugar ‘stingy \dots\ mood’}\Bfootnote{Ettins are characteristically inhospitable, in stark opposition to the Old Germanic social norms; see Introduction to the poem above.  This statement foreshadows the later hunting expedition starting at st.~16 below.}}.“\eva

\bvb\speakernoteb{[Tew’s mother:]}“O clansman of ettins \ken*{= Tew}! I would wish to place \\
you two, full of heart, beneath the cauldrons. \\
Many a time has my lover \ken*{= Hymer} been \\
stingy with guests, quick to ill mood.”\evb\evg


\bvg\bva\mssnote{\Regius~14r/11, \AM~6r/9}%
En \alst{v}á-skapaðr \hld\ \alst{v}arð \edtrans{síð-búinn}{delayed}{\Afootnote{om. \AM}}, &
\alst{h}arð-ráðr \alst{H}ymir, \hld\ \alst{h}ęim af vęiðum; &
\alst{g}ekk inn ï sal, \hld\ \alst{g}lumðu \edtrans{jǫklar}{icicles}{\Bfootnote{In Hymer’s frozen beard.  In modern Icelandic the word \emph{jökull} has come to mean ‘glacier’, but its original sense (as found here) is that of its English cognate “icicle”.}}, &
vas \edtrans{\alst{k}arls}{churl}{\Bfootnote{The first of three instances in the poem of Hymer being called a \emph{karl} ‘churl, old man, peasant’; a reference to his boorish nature.  For the sense of this word in ON poetry (in contrast to both ON prose and later Scandinavian languages) see \textlink{Harbardsljod}[2] n.}}, es \alst{k}om, \hld\ \alst{k}inn-skógr frørinn.\eva

\bvb But the misshapen one was delayed, \\
hard-minded Hymer, home from the hunt. \\
He entered the hall; icicles clattered; \\
on the churl who came was the cheek-shaw \ken{beard} frozen.\evb\evg


\bvg\bva\speakernote{[Týs móðir:]}\mssnote{\Regius~14r/13, \AM~6r/11}%
\llap{„}\edtext{Ves \alst{h}ęill, \alst{H}ymir, \hld\ ï \alst{h}ugum góðum!}{\lemma{Ves hęill, \dots\ ï hugum góðum! ‘Be thou hale \dots\ in good spirits!’}\Bfootnote{A formulaic greeting; cf.~the almost identical greeting in N B380 (edited below under Galders).  Further afield cf.~the type exemplified by \Beowulf\ 407a: \emph{Wæs þú, Hróð·gâr, hâl} ‘Be thou, Rothgar, hale!’}} &
Nú ’s \alst{s}onr kominn \hld\ til \alst{s}ala þïnna, &
sá’s \alst{v}it \alst{v}ę̇ttum \hld\ af \alst{v}egi lǫngum; &
fylgir \alst{h}ǫ̇num \hld\ \alst{H}róðrs and·skoti, &
\alst{v}inr \alst{v}er-liða; \hld\ \edtrans{\alst{V}é·urr}{Wighward}{\Bfootnote{The guardian of \inx[C]{wigh}[wighs] (sanctuaries), a name of Thunder.}} hęitir sá.\eva

\bvb\speakernoteb{[Tew’s mother:]}“Be thou hale, Hymer, in good spirits! \\
Now the son has come to thy halls, \\
he whom we’ve awaited, from a long way off. \\
He is followed by Rooder’s enemy \ken*{= Thunder}, \\
the friend of manly retinues---\inx[P]{Wighward} is he called.\evb\evg


\bvg\bva\mssnote{\Regius~14r/15, \AM~6r/13}%
\alst{S}é hvar \alst{s}itja \hld\ und \alst{s}alar gafli, &
\alst{s}vá \edtext{forða \alst{s}ér}{\Afootnote{\emph{forðask} \AM}}, \hld\ stęndr \edtrans{\alst{s}úl}{column}{\Afootnote{\emph{†sol†} \AM}} fyrir.“ &
\alst{S}undr stǫkk \alst{s}úla \hld\ fyr \alst{s}jȯn jǫtuns, &
ęn \alst{á}ðr ï tvau \hld\ \alst{ȧ}ss brotnaði.\eva

\bvb See where they sit beneath the hall’s gable: \\
so they save themselves---a column stands before them!” \\
The column burst apart before the ettin’s gaze, \\
but earlier the roof-beam broke in two.\evb\evg


\bvg\bva\mssnote{\Regius~14r/17, \AM~6r/15}%
\edtext{Stukku \alst{á}tta, \hld\ ęn \alst{ęi}nn af þęim &
\alst{h}verr \alst{h}arð-slęginn \hld\ \alst{h}ęill af þolli;}{\lemma{Stukku \dots\ þolli ‘Eight \dots\ peg.’}\Bfootnote{Nine cauldrons were hanging from the roof-beam supported by the column behind which the gods were hiding.  Eight of them broke from the fall and one remained whole.  We may presume that this was the famed cauldron the gods had come to get.}} &
\alst{f}ramm gingu þęir, \hld\ ęn \alst{f}orn jǫtunn &
\alst{s}jȯnum lęiddi \hld\ \alst{s}ïnn and·skota.\eva

\bvb Eight [cauldrons] crashed down, but one of them, \\
a hard-forged cauldron, [came] whole off its peg. \\
Forth they went, but the ancient ettin \\
with his gaze tracked his enemy.\evb\evg


\bvg\bva\mssnote{\Regius~14r/19, \AM~6r/16}%
\edtrans{Sagði⸗t \alst{h}ǫ̇num \hld\ \alst{h}ugr vel}{His heart did not please him}{\Bfootnote{Lit. “his heart did not speak well to him”.}} þȧ’s sá &
\alst{g}ýgjar \edtrans{\alst{g}rǿti}{distresser}{\Afootnote{\emph{gę́ti} ‘keeper, warder’ \AM}} \hld\ ȧ \alst{g}olf kominn, &
\alst{þ}ar vǫ́ru \alst{þ}jórar \hld\ \alst{þ}rír of tęknir, &
bað \edtrans{\alst{s}ęnn}{at once}{\Afootnote{\emph{sun} ‘[his] son \ken*{= Tew}?’ \AM}} jǫtunn \hld\ \alst{s}jóða ganga.\eva

\bvb His heart did not please him when he saw \\
the \inx[C]{gow}’s distresser \ken*{= Thunder} come upon the floor. \\
There were three bulls a-taken; \\
the ettin bade them at once go cook.\evb\evg


\bvg\bva\mssnote{\Regius~14r/21, \AM~6r/18}%
\alst{H}vęrn létu þęir \hld\ \alst{h}ǫfði skęmmra &
auk ȧ \alst{s}ęyði \hld\ \alst{s}íðan bǫ́ru, &
\edtext{át \alst{S}ifjar verr \hld\ áðr \alst{s}ofa gingi, &
\alst{ęi}nn með \alst{ǫ}llu \hld\ \alst{ø}xn tvȧ Hymis.}{\lemma{át \dots\ Hymis. ‘Sib’s \dots\ oxen.’}\Bfootnote{Cf.~\textlink{Thrymskvida}[24] for another instance of Thunder’s great eating, which curiously also uses the kenning \emph{Sifjar verr} ‘Sib’s husband \ken*{= Thunder}’.}}\eva

\bvb Each one they made shorter by a head \\
and onto the cooking-pit then did carry. \\
Sib’s husband \ken*{= Thunder} ate, before he went to sleep, \\
alone by himself two of Hymer’s oxen.\evb\evg


\bvg\bva\mssnote{\Regius~14r/23, \AM~6r/19}%
\Ballnote{Hymer’s stinginess---he refuses to share more of his own food but instead forces his guests to go hunt---breaks all Indo-European rules of hospitality and illustrates the otherness of the Ettins; see Introduction above.}%
Þȯtti \alst{h}ǫ́rum \hld\ \alst{H}rungnis spjalla &
\alst{v}erðr Hlór·riða \hld\ \alst{v}ęl full-mikill, &
\speakernote{[Hymir:]}%
\llap{„}\edtext{munum at \alst{a}ptni \hld\ \alst{ǫ}ðrum verða &
\alst{v}ið \alst{v}ęiði-mat \hld\ \alst{v}ér þrír lifa.“}{\lemma{munum \dots\ lifa. ‘Next \dots\ live.’}\Bfootnote{The poet is pushing at the limits of Old Norse syntax.  In prose word order it should be construed as: \emph{at ǫðrum aptni munum vér þrír verða lifa við vęiði-mat}, where \emph{verða} ‘have to, must’ is used like its modern German cognate \emph{werden}.}}\eva

\bvb To Rungner’s hoary friend \ken*{= Hymer} did seem \\
Loride’s eating far too great--- \\
“Next evening will we three \\
on game-meat have to live.”\evb\evg


\bvg\bva\mssnote{\Regius~14r/24, \AM~6r/21}%
\alst{V}é·urr kvaðsk \alst{v}ilja \hld\ ȧ \alst{v}ág róa, &
ef \alst{b}allr jǫtunn \hld\ \alst{b}ęitur gę́fi. &
\speakernote{[Hymir:]}%
\llap{„}\alst{H}verf þú til \edtext{\alst{h}jarðar}{\Afootnote{\emph{hallar} corr. \AM}}, \hld\ ef \alst{h}ug trúir, &
\edtrans{\alst{b}rjótr \alst{b}erg-Dana}{breaker of boulder-Danes \ken*{\textsc{ettins} > = Thunder}}{\Bfootnote{This kenning for Thunder also occurs in \Haustlong\ 18; it emphasises the otherness of the Ettins (see Introduction to the poem) by equating them with ethnic foreigners.  Cf.~also \Thorsdrapa, where ettins are called Scots, Swedes, Danes, \inx[P]{Ruges} and \inx[P]{Harths}, all peoples hostile to the Norwegian Earl Hathkin, at whose court that poem may have been composed.}}, \hld\ \alst{b}ęitur sǿkja.\eva

\bvb Wighward called himself willing to row on the wave \\
if the stubborn ettin might give pieces of bait. \\
\speakernoteb{[Hymer:]}%
“Turn to the herd if thou trustest thy heart--- \\
O breaker of boulder-Danes \ken*{\textsc{ettins} > = Thunder}---to seek pieces of bait.\evb\evg


\bvg\bva\mssnote{\Regius~14r/26, \AM~6r/23}%
\alst{Þ}ęss \edtext{vę̇ntir mik}{\Afootnote{so \AM; \emph{vę̇nti ek} (norm.) \Regius}}, \hld\ at \alst{þ}ér \edtrans{myni⸗t}{will not}{\Afootnote{so \AM; \emph{myni} ‘will’ \Regius.  The \AM\ reading is preferable since it makes this the first of Hymer’s several challenges of strength to Thunder, which the god, to the ettin’s humiliation, easily accomplishes.}} &
\alst{ǫ}gn at \alst{o}xa \hld\ \alst{au}ð-fęng vesa.“ &
\edtrans{\alst{S}vęinn}{The swain}{\Bfootnote{Thunder was in the shape of a young boy, a detail which is also found in \Gylfaginning~48:4; see Introduction to the poem.}} \alst{s}ýsliga \hld\ \alst{s}vęif til skógar, &
þar’s \edtext{\alst{o}xi stóð \hld\ \alst{a}l-svartr}{\lemma{oxi \dots\ al-svartr ‘ox \dots\ all-black’}\Bfootnote{Formulaic, also occuring in \textlink{Thrymskvida}[23]; see note there for further parallels to the custom of sacrificing animals of certain colours.  It seems that all-black oxen were thought the noblest, and so Thunder’s slaying one instead of an inferior beast is probably intended to humiliate the stingy Hymer. ---
We learn more about this ox in \Gylfaginning~48:13: \emph{Hann tók inn mesta uxa’nn, er Himin-hrjóðr hét, ok sleit af hǫfuð’it ok fór með til sjávar.} ‘He took the greatest ox, which was called Heavenrid, and tore off its head and went with it to the sea’.}} fyrir.\eva

\bvb I expect that the bait from the ox \\
will not be an easy catch for thee!” \\
The swain \ken*{= Thunder} swiftly turned to the wood \\
where an ox stood, all-black, ahead.\evb\evg


\bvg\bva\mssnote{\Regius~14r/28, \AM~6r/24}%
Braut af \alst{þ}jóri \hld\ \alst{þ}urs ráð-bani &
\alst{h}ǫ́-tu̇n ofan \hld\ \alst{h}orna tvęggja. &
\speakernote{[Hymir:]}%
\llap{„}\alst{V}erk þikkja þïn \hld\ \alst{v}erri myklu &
\edtrans{\alst{k}jóla valdi}{wielder of keels}{\Bfootnote{I.e.~‘master of ships’, foreshadowing the later fishing expedition. --- \emph{valdi} is here taken as dat.~sg.~of \emph{valdr} ‘wielder, master’ and as the indirect object of \emph{þykkja} ‘seem, appear’.  The word can however also be read as the semantically identical nom.~sg.~\emph{valdi}, in which case the phrase would be a vocative addressing Thunder.
In context the former reading makes more sense: the kenning ‘ruler of keels’ fits better for Hymer than Thunder and \emph{þykkja} gets an indirect object.}} \hld\ an \alst{k}yrr sitir.“\eva

\bvb Off from the bull the thurse’s death-planner \ken*{= Thunder} \\
broke the high meadow of two horns \ken{head} from above.--- \\
\speakernoteb{[Hymer:]}%
“Worse by far thy works appear \\
to the wielder of keels \ken*{= Hymer = me} than if thou wouldst sit calm!”\evb\evg


\bvg\bva\mssnote{\Regius~14r/30, \AM~6r/26}%
\Ballnote{The scene shifts and the party is out at sea.  It is possible that a stanza has been lost here, or that the transition would have been indicated in some other way during oral performance.}%
Bað \alst{h}lunn-gota \hld\ \edtrans{\alst{h}afra dróttinn}{Lord of He-goats}{\Bfootnote{Thunder.  Reoccurs in 31/1 below.}} &
\edtext{\alst{á}tt-runn}{\Afootnote{\emph{†atrænn†} \AM}} \edtrans{\alst{a}pa}{ape}{\Bfootnote{The specific sense of \emph{api} ‘ape’ is uncertain.  It seems to generally refer to a fool, but see Index.}} \hld\ \alst{ú}tar fǿra, &
\edtext{en \alst{s}á jǫtunn \hld\ \alst{s}ïna \edtext{talði}{\Afootnote{\emph{milldi} corr. \AM}}, &
\alst{l}ítla fẏsi \hld\ \edtext{at róa \alst{l}ęngra}{\Bfootnote{The alliteration falls on the second stressed word which is not expected; emendation to \emph{*lęngra at róa} is metrically convenient but not supported by either ms.}}}{\lemma{en \dots\ róa. ‘but \dots\ longer.’}\Bfootnote{Thunder’s humorous humiliation of Hymer continues with the snide ettin now forced to row against his will.}}.\eva

\bvb The Lord of He-goats bade the kinsman of the \inx[C]{ape}\ \ken*{\textsc{ettin} = Hymer} \\
push the launcher-steed \ken{boat} further out, \\
but that ettin told of his \\
little wish to row any longer.\evb\evg


\bvg\bva\mssnote{\Regius~14r/31, \AM~6r/27}%
Dró \edtrans{\alst{m}ę́rr}{famous}{\Afootnote{so \Regius; \emph{męirr} ‘more, further’ \AM}} Hymir \hld\ \alst{m}óðugr \edtext{hvali}{\Bfootnote{A rare acc. pl. form also occurring in 26/2b.}} &
\alst{ęi}nn ȧ \alst{ǫ}ngli \hld\ \alst{u}pp sęnn tvȧ\emph{a}; &
en \alst{a}ptr ï skut \hld\ \alst{Ó}ðni sifjaðr &
\alst{V}é·urr við \alst{v}élar \hld\ \alst{v}að gęrði sér.\eva

\bvb Famous Hymer, fierce, pulled whales: \\
one on the hook, soon up two--- \\
but back in the stern Weden’s kin, \\
Wighward \name{= Thunder} craftily fixed his line.\evb\evg


\bvg\bva\mssnote{\Regius~14v/1, \AM~6r/29}%
\Ballnote{This stanza skilfully juxtaposes Thunder and the Wyrm by employing strong parallellism across the two helmings (ll. 1--2 and 3--4).  The poetic juxtaposition between the two primordial enemies, Storm-God and Serpent, is also found in the Scaldic fragments (e.g. Ǫlv \emph{Þórr}; see note to l.~4) and, further afield, in the Vedic tradition (e.g. \Rigveda\ 1.32.13c: \textoi{Índraş ca yád \hld\ yuyudʰā́tē Áɦiş ca} ‘When Índra and the Wyrm (\textoi{áɦi}) fought each other.’)}%
\alst{Ę}gnði ȧ \alst{ǫ}ngul \hld\ sá’s \alst{ǫ}ldum bergr, &
\edtrans{\alst{o}rms \alst{ęi}n-bani}{Wyrm’s lone slayer}{\Bfootnote{\textcite[419--424]{Watkins1995} sees in this kenning for Thunder a reflex of the ancient IE formula \textsc{hero slay} (PIE \emph{*gʷʰen-}) \textsc{serpent},
where the root \emph{*gʷʰen-} ‘slay, kill’ is reflected in ON by the agent noun \emph{bani} ‘bane, slayer’ < \emph{*gʷʰon-ō} (so e.g. \cite[113--114]{Seebold1967}, although the sound change \emph{*gʷʰ-} > \emph{*b-} is not universally accepted; see \cite[xxxviii,51]{Kroonen2013}).
The IE formula, reconstructed as \emph{*h₁é·gʷʰent h₁ógʷʰim} ‘he slew the Serpent’, occurs directly only as OI \textoi{áɦann Áɦim} (11 times in \Rigveda; see \cite[304\psqq]{Watkins1995}), but has parallels throughout the IE family, including Greek and Hittite \parencite[\textit{passim}]{Watkins1995}.

On the basis of the rare sound combination \emph{*-gʷʰ- : -gʷʰ-}, \textcite{Wikander2021} argues brilliantly that the IE formula was likely phono-semantically calqued from IE into Northwest Semitic as the formulaic rhyming phrase \emph{*maḫaṣ́a naḥaša(m)} ‘he struck the serpent’ (which features another rare and highly marked phoneme, the ejective \emph{*ṣ́} \textipa{[\super t\textbeltl ']}), perhaps along with the Chaoskampf myth as a whole.}} \hld\ \alst{o}xa hǫfði; &
\alst{g}ęin við \edtrans{agni}{bait}{\Afootnote{so \AM; \emph{ǫngli} ‘hook’ \Regius}} \hld\ sú’s \alst{g}oð fía &
\edtext{\alst{u}mb-gjǫrð neðan \hld\ \alst{a}llra landa.}{\lemma{umb-gjǫrð \dots\ allra landa ‘engirdler of all lands’}\Bfootnote{The same kenning is found in a fragment by Alewigh Snub (Ǫlv \emph{Þórr} in \Skp\ 3) quoted in \Skaldskaparmal\ 11: \emph{\alst{Ǿ}sti⸗sk \alst{a}llra landa \hld\ \alst{u}mb-gjǫrð ok sonr Jarðar} ‘The engirdler of all lands and the son of Earth surged.’
Closely related is the kenning in a 9th c.~fragment by Braye the Old quoted in the same chapter (Bragi \emph{Þórr} 3 in \Skp\ 3): \emph{ęndi-sęiðr allra landa} ‘boundary-saithe of all lands’.  Both kennings relate to the cosmological idea of the Outer Sea (ON \emph{Út-haf}) which surrounds the solid Earth; in this sea lies the Wyrm, biting its tail and enwrapping the lands.  Cf.~st.~5/1 n.~above.}}\eva

\bvb On the hook baited he who rescues men---  \\
the Wyrm’s lone slayer---the ox’s head. \\
At the bait snapped the one whom the Gods hate--- \\
the engirdler of all lands---from below.\evb\evg


\bvg\bva\mssnote{\Regius~14v/3, \AM~6v/1}%
\alst{D}ró \alst{d}jarf⸗liga \hld\ \alst{d}áð-rakkr \edtrans{Þȯ\emph{a}rr}{Thunder}{\Bfootnote{Out of 8 three-syllable lines in \textlink{Hymiskvida}, this is the only one which is deficient in both \Regius\ and \AM\ and cannot easily be emended by restoring an hiatus form.  In the quite strict meter (see Introduction above) observed by the \textlink{Hymiskvida} poet we should expect a disyllabic form in this spot, and this may be had if we restore an archaic \emph{*Þȯarr} or \emph{*Þȯurr}.  Such a form is less secure than other hiatus forms, but is also metrically required in \textlink{Hymiskvida}[28]/2b below and further in the Scaldic \Thorsdrapa\ 2/2b.
This issue is treated in depth by \textcite{Haukur2023}, who argues for the form \emph{*Þóurr}.}} &
\alst{o}rm \alst{ęi}tr-fá\emph{a}n \hld\ \alst{u}pp at borði; &
\alst{h}amri kníði \hld\ \edtrans{\alst{h}ǫ́-fjall skarar}{high mountain of hair \ken{head}}{\Bfootnote{A rather unfitting kenning, since serpents do not have hair.}} &
\alst{o}f-ljótt \alst{o}fan \hld\ \alst{u}lfs hnit-bróður.\eva

\bvb Bravely deed-ready Thunder pulled \\
the venom-gleaming Wyrm up on the gunwale. \\
With the hammer he struck the high mountain of hair \ken{head}--- \\
very hideous, from above---on the Wolf’s clash-brother \ken*{= the Wyrm}.\evb\evg


\bvg\bva\mssnote{\Regius~14v/5, \AM~6v/2}%
\edtrans{\alst{H}r\emph{au}n-gǫlkn}{The waste-monsters}{\Afootnote{emend.; \emph{hręin-gǫlkn} \AM\Regius}\Bfootnote{Both mss. have \emph{hręin-}, which may mean either ‘clean’ or ‘reindeer’, neither of which make sense Semantically.  On the other hand \emph{hraun} \ONP: ‘stone/barren area, wasteland; lavafield’ is well attested in Scaldic kennings for ettins, including in the present poem (36/4, 38/3).  The precise meaning of \emph{galkn} ‘monster’ (plural \emph{gǫlkn}) is unclear, but it is attested in 3 Scaldic verses, in all cases as the base-word in kennings of the type “troll-woman of the shield \ken{axe}”, which may suggest that it (like another neuter word, \emph{flagð}) refers specifically to female malevolent beings.

While the mss. spelling ‘\emph{galkn}’ (norm. \emph{gálkn}) could reflect either singular or plural, the form of the verb is plural.  This means that the word cannot be referring to the Middenyardswyrm as \textcite{LarringtonEdda} understands it in her translation: “the sea-wolf shrieked”.}} \edtext{\alst{h}rutu}{\lemma{hrutu}\Afootnote{so \AM; \emph{hlumðu} ‘dashed’ \Regius.}\Bfootnote{The \AM\ reading is preferred since it has the metrically required short root syllable.  End rhyme is parallelled elsewhere in the poem (st.~3/3); here it heightens the dramatic impact of the verse.}}, \hld\ en \alst{h}ǫlkn þutu, &
\alst{f}ór hin \alst{f}orna \hld\ \alst{f}old ǫll saman; &
\skipnumbering \edtext{[\dots]}{\Bfootnote{It is very likely that a line is missing here, since the stanzas in the poem fairly consistently have 4 lines and since the \emph{síðan} ‘thereafter’ in l.~3 is left without a referent.  In other tellings of the myth it is at this point that Hymer cuts Thunder’s fishing line, so that is probably what has been lost.  Cf.~\Gylfaginning~48:26--27: \emph{Þá er sagt, at jǫtunn’inn Hymir gerði⸗sk lit-verpr, fǫlnaði ok hrę́ddi⸗sk, er hann sá orm’inn ok þat er sę́r’inn féll út ok inn of nǫkkvann.  Ok í því bili, er Þórr greip hamar’inn ok fǿrði á loft, þá fálmaði jǫtunn’inn til agn-saxi’nu ok hjó vað Þórs á borði, en ormr’inn søkkði⸗sk í sę́’inn.} ‘Then it is said that the ettin Hymer lost his colour, turned pale and was frightened when he saw the Wyrm and how the sea fell in and out of the boat.  And in that moment when Thunder gripped the hammer and brought it aloft then the ettin reached for the bait-cutting knife and cut Thunder’s fishing-line against the gunwale.  But the Wyrm sunk into the sea.’

Based on this account we may reconstruct a hypothetical line \emph{*und’s fǫlr Hymir \hld\ fekk ȧ saxi} ‘until pale Hymer grasped the knife’.}} &
\alst{s}økkði⸗sk \alst{s}íðan \hld\ \alst{s}á \edtrans{fiskr}{fish}{\Bfootnote{A good example showing that pre-modern speciation, especially in poetry, was not exact.  The appellation is not a mere fancy of the poet, for the Wyrm appears in the shape of a fish in various pre-Christian pictorial sources (e.g. GP 21).  It may also be called a fish in \textlink{Grimnismal}[21] (see note there), and in Scaldic sources it is often called a saithe (\emph{sęiðr}).}} ï mar.\eva

\bvb The waste-monsters \ken{ettins} bounded and bedrock resounded; \\
the ancient earth moved all at once. \\
{[\dots]} \\
Thereafter sank that fish \ken*{= the Wyrm} into the sea.\evb\evg


\bvg\bva\mssnote{\Regius~14v/6, \AM~6v/3}%
\alst{Ȯ}·tęitr \alst{jǫ}tunn, \hld\ es \alst{a}ptr røru, &
\skipnumbering \edtext{[\dots]}{\Bfootnote{Another likely missing line.  As noted in the previous stanza the meter usually requires four lines; more importantly the first half of the sentence is incomplete without a verb.}} &
\edtrans{svá’t \alst{á}r}{so that for a long time}{\Bfootnote{\textcite{FinnurEdda}\ suggests \emph{svá’t at ǫ́r} ‘so that by the oar’, but this burdens the strict meter.  For this sense of \emph{ár} cf.~\textlink{Skirnismal}[27].}} Hymir \hld\ \alst{ę}kki mę́lti, &
\alst{v}ęifði rǿði \hld\ \alst{v}eðrs annars til.\eva

\bvb The unmerry ettin \ken*{= Hymer} as they rowed back, \\
{[\dots]} \\
so that for a long time Hymer said not a thing; \\
he pulled the oar against the wind:\evb\evg


\bvg\bva\speakernote{[Hymir:]}\mssnote{\Regius~14v/8, \AM~6v/4}%
\Ballnote{Hymer tells Thunder, who having let go of the Wyrm has nothing to show for the trip, that he can accomplish something half as great as the catching of the whales if he carries them home and ties the boat by the shore.}%
\llap{„}Munt of \alst{v}inna \hld\ \alst{v}erk halft við mik, &
at \alst{h}ęim \alst{h}vali \hld\ \alst{h}af til bǿjar &
eða \alst{f}lot-brúsa \hld\ \alst{f}ęstir okkarn.“\eva

\bvb\speakernoteb{[Hymer quoth:]}%
“Thou wilt accomplish a half work compared to me, \\
if thou bringest home the whales to the farm, \\
or our float-buck \ken{boat} thou dost fasten.”\evb\evg


\bvg\bva\mssnote{\Regius~14v/9, \AM~6v/6}%
\alst{G}ekk Hlór·riði \hld\ \alst{g}ręip \edtext{ȧ}{\Afootnote{\emph{til ȧ} \Regius}} stafni &
vatt \edtrans{með \alst{au}stri}{with the bilge-water}{\Bfootnote{That is, the bilge-water was still inside the boat; another impressive work of strength.}} \hld\ \alst{u}pp lǫg-fáki; &
\alst{ęi}nn með \alst{ǫ́}rum \hld\ ok með \alst{au}st-skotu &
\alst{b}ar til \alst{b}ǿjar \hld\ \alst{b}rim-svïn jǫtuns &
ok \edtext{\edtext{\alst{h}olt-riða}{\Afootnote{\emph{†holtriba†} \Regius}} \hld\ \alst{h}ver}{\lemma{holt-riða hver ‘spring of woodland ridges’}\Bfootnote{An uncertain geographical description apparently meant to impress the original audience.  In Iceland \emph{hverr} ‘cauldron’ also carries the sense ‘(hot) spring’, in which case this could be an attestation of Thunder’s prowess in wading (for which see \textlink{Grimnismal}[29]).  TODO: What do other editors and translators say?}} ï gęgnum. \eva

\bvb Loride went, grasped the stern, \\
hurled up the lake-nag \ken{boat} with the bilge-water. \\
Alone with the oars and the bilge-bucket \\
he bore to the farm the ettin’s brim-swines \ken{whales}, \\
even through the spring of woodland ridges.\evb\evg


\bvg\bva\mssnote{\Regius~14v/12, \AM~6v/7}%
\Ballnote{Even after witnessing several great feats of strength Hymer still refuses to admit Thunder’s superiority.  He now insists on challenging him to break his indestructible chalice.}%
\edtext{Ok}{\Afootnote{\emph{Enn} \AM}} \alst{ę}nn \alst{jǫ}tunn \hld\ umb \alst{a}fr-endi, &
\alst{þ}rá-girni vanr, \hld\ við \alst{Þ}ȯ\emph{a}r sęnti, &
kvað⸗at mann \alst{r}amman, \hld\ þó’tt \alst{r}óa kynni, &
\alst{k}rǫptur⸗ligan, \hld\ nema \alst{k}alk bryti.\eva

\bvb And yet the ettin, used to stubbornness, \\
over strength of hand with Thunder flyted. \\
He would call no man strong, even though he could row \\
mightily, unless he broke the chalice.\evb\evg


\bvg\bva\mssnote{\Regius~14v/14, \AM~6v/9}%
En \alst{H}lór·riði, \hld\ es at \alst{h}ǫndum kom, &
\alst{b}rátt lét \alst{b}resta \hld\ \edtrans{\alst{b}ratt-stęin glęri}{steep stone with the glass}{\Bfootnote{Thunder broke the stone columns in Hymer’s house with the glass chalice.}}, &
\alst{s}ló \edtrans{\alst{s}itjandi}{standing}{\Bfootnote{This word is ambiguous and can modify either Thunder (in which case it would mean “sitting”) or the columns (\emph{súlur}).  I have chosen the latter and read it as signifying their stability.}} \hld\ \alst{s}úlur ï gǫgnum; &
bǫ́ru þó \alst{h}ęilan \hld\ fyr \alst{H}ymi síðan,\eva

\bvb But Loride, when it came to his hands, \\
impatiently crushed steep stone with the glass. \\
He struck right through the standing columns, \\
yet they brought it, whole, before Hymer thereafter,\evb\evg


\bvg\bva\mssnote{\Regius~14v/16, \AM~6v/10}%
\edtext{und’s þat hin \alst{f}ríða \hld\ \alst{f}riðla kęndi &
\alst{ǫ̇}st-ráð mikit, \hld\ \alst{ęi}tt es vissi,}{\lemma{und’s \dots\ vissi, ‘until \dots\ knew:’}\Bfootnote{Harkening back to st.~4.}} &
\llap{„}drep við \edtrans{\alst{h}aus \alst{H}ymis, \hld\ hann ’s \alst{h}arðari}{Hymer’s skull! It is harder}{\Bfootnote{Ettins characteristically have hard skulls; cf.~\textlink{Harbardsljod}[15] and note for the significance of this minor detail.}}, &
\edtrans{\alst{k}ost-móðs}{choice-weary}{\Bfootnote{Hymer can justly be called “choice-weary” since at this point the gods have destroyed eight of his nine cauldrons and slain three of his bulls.}} jǫtuns, \hld\ \alst{k}alki hvęrjum!“\eva

\bvb until the handsome mistress \ken*{= Tew’s mother} gave \\
a great loving counsel, the one she knew: \\
“Strike against Hymer’s skull! It is harder--- \\
the choice-weary ettin’s---than any chalice.”\evb\evg


\bvg\bva\mssnote{\Regius~14v/18, \AM~6v/12}%
\alst{H}arðr \edtext{ręis}{\Afootnote{om. \AM}} ȧ kné \hld\ \alst{h}afra dróttinn, &
\edtrans{fǿrðisk \alst{a}llra \hld\ ï \alst{ǫ̇}s-męgin}{brought himself into his fullest Os-might}{\Bfootnote{The exact meaning of this is not clear, but a likely reading is that Thunder brings himself into his true form as the mighty, red-bearded man; recall that he was previously (st.~18/3) in the shape of a young boy.  Cf.~\Gylfaginning~48:24 in its description of Thunder attempting to pull up the Wyrm: \emph{Þá varð Þórr reiðr ok fǿrði⸗sk í ǫ́s-megin [\dots]} ‘Then Thunder turned wroth and brought himself into his Os-might’ and Eddic fragment 6, st.~1.}}; &
\alst{h}ęill vas karli \hld\ \alst{h}jalm-stofn ofan, &
en \alst{v}ïn-fęrill \hld\ \alst{v}alr rifnaði.\eva

\bvb On his knee rose the hard Lord of He-goats; \\
he brought himself to his highest Os-might. \\
Whole on the churl \ken*{= Hymer} was the helm-stump \ken{head} above, \\
but the round wine-track \ken{chalice} rent apart.\evb\evg


\bvg\bva\speakernote{[Hymir kvað:]}\mssnote{\Regius~14v/20, \AM~6v/13}%
\llap{„}\alst{M}ǫrg vęit’k \alst{m}ę́ti \hld\ \alst{m}ér gingin frȧ, &
\edtext{es}{\Afootnote{om. \Regius}} \alst{k}alki sé’k \hld\ \edtext{fyr}{\Afootnote{\emph{ýr} \Regius; \emph{firi} \AM}} \edtext{\alst{k}n\emph{éu}m}{\Afootnote{metr. emend.; \emph{knjám} (norm.) \Regius\AM}} hrundit,“ &
\alst{k}arl orð of \alst{k}vað: \hld\ „\edtext{\alst{k}ná’k⸗at sęgja &
\alst{a}ptr \alst{ę́}va-gi: \hld\ ‚þú ’st \alst{ǫ}lðr of hęitt.}{\lemma{kná’k⸗at \dots\ of hęitt. ‘I cannot \dots\ warmed!’}\Bfootnote{Hymer laments that with the loss of his finest vessel he will never be able to enjoy his drink again.  This is ironic since it was he who challenged Thunder to break it in the first place.}}‘\eva

\bvb\speakernoteb{[Hymer quoth:]}%
“I know many treasures parted from me \\
when I see the chalice thrown before [my] knees!”--- \\
The churl \ken*{= Hymer} spoke words: “I cannot say \\
ever again: ‘Thou art, ale, well warmed!’\evb\evg


\bvg\bva\mssnote{\Regius~14v/22, \AM~6v/15}%
Þat ’s til \alst{k}ostar \hld\ ef \alst{k}oma mę́ttið &
\alst{ú}t ór \alst{ȯ}ru \hld\ \edtrans{\alst{ǫ}l-kjól}{ale-vessel \ken{cauldron}}{\Bfootnote{\emph{ǫl-kjól} is the accusative of \emph{ǫl-kjóll}, but in this construction (\CV: \emph{koma}, B) we would expect the dative \emph{ǫl-kjóli}.  Since the meter does not allow for this the poet has probably taken a grammatical liberty.}} \edtrans{hofi}{hall}{\Bfootnote{This is the only Old Norse occurrence of the word \emph{hof} in the sense “hall, house”---it otherwise only means “temple” (\inx[C]{hove}).  The West Germanic cognates consistently mean “hall”, but that is probably the original sense, so it is unclear if this is an instance of foreign (if so, most likely Anglo-Saxon) influence or just a poetic archaism.}}.“ &
\alst{T}ýr lęitaði \hld\ \alst{t}ysvar hrǿra; &
stóð at \alst{h}vǫ́ru \hld\ \alst{h}verr kyrr fyrir.\eva

\bvb It would be choicest if ye might bring \\
that ale-vessel \ken{cauldron} out from our hall.” \\
Tew attempted twice to move it--- \\
each time stood the cauldron unmoved ahead.\evb\evg


\bvg\bva\mssnote{\Regius~14v/24, \AM~6v/16}%
\alst{F}aðir Móða \hld\ \alst{f}ekk ȧ þręmi &
\edtrans{ok ï \alst{g}ęgnum \edtrans{stęig}{stepped}{\Afootnote{so \Regius; \emph{stóð} ‘stood’ \AM}} \hld\ \alst{g}olf niðr \edtext{ï}{\Afootnote{\emph{ȧ} (norm.) \AM}} sal}{and stepped through the floor in the hall}{\Bfootnote{In \Gylfaginning~48:24 Thunder is said to have stepped through the boat when trying to pull up the Middenyardswyrm, a detail famously depicted on the Altuna runestone (U 1611) from Uppland, Sweden and the Hørdum stone from Denmark.  It is very likely that the original detail of Thunder exerting himself to the point of his feet going through the bottom of the boat has later been transposed to this place in the narrative, although it is unclear whether this is the work of the poet himself or of an earlier tradition.}}; &
hóf⸗sk ȧ \edtrans{\alst{h}\emph{au}fuð}{head}{\Bfootnote{The archaic long form \emph{haufuð} rather than the short \emph{hǫfuð} is metrically superior in a type B verse (TODO: cite Males2025 note).  Cf.~the similar line 36/1 below.}} upp \hld\ \alst{h}ver Sifjar verr, &
en ȧ \alst{h}ę̇lum \hld\ \edtrans{\alst{h}ringar skullu}{the rings clattered}{\Bfootnote{I.e. the chain-links.  This detail is mentioned in an example sentence contrasting long and short phonemes in \FGT: \emph{heyrði til hǫddu, þá er Þórr bar hverinn} ‘the sound of the pot-links (\emph{hadda}) was heard when Thunder bore the cauldron’.  According to \textcite{FinnurEdda} the chain (or \emph{hadda}) on a Wiking-age cauldron would reach across it, being bound on two opposing points along the rim, in which case this would be a reference to the cauldron’s enormous size, with its diameter---mentioned in st.~5 as one \inx[C]{rest}---being roughly the same as Thunder’s height.}}.\eva

\bvb The father of Moody \ken*{= Thunder} grasped the brim \\
and stepped through the floor in the hall. \\
Up onto his head heaved Sib’s husband \ken*{= Thunder} the cauldron, \\
but by his heels the rings clattered.\evb\evg


\bvg\bva\mssnote{\Regius~14v/26, \AM~6v/18}%
Fóru⸗t \alst{l}ęngi, \hld\ áðr \alst{l}íta nam &
\alst{a}ptr \alst{Ó}ðins sonr \hld\ \alst{ęi}nu sinni; &
sá ór \alst{h}ręysum \hld\ með \alst{H}ymi austan &
\edtext{\alst{f}olk-drótt}{\lemma{folk-drótt \dots\ fjǫl-hǫfðaða ‘war-troop \dots\ many-headed’}\Bfootnote{The adjective \emph{fjǫl-hǫfðaðr} means ‘many-headed, polycephalic’ and is not referring to the size of the host.  For many-headed ettins see st.~8 and for their malformed bodies in general see Introduction to the poem.}} \alst{f}ara \hld\ \alst{f}jǫl-hǫfðaða.\eva

\bvb They did not journey for long before Weden’s son \\
took to look back a single time. \\
He saw out of stone-heaps with Hymer from the east \\
a war-troop coming, many-headed.\evb\evg


\bvg\bva\mssnote{\Regius~14v/28, \AM~6v/19}%
\Ballnote{For another light-hearted instance of Thunder’s ettin-bashing cf.~\textlink{Thrymskvida}[31]--32.}%
\alst{H}óf sér af \alst{h}ęrðum \hld\ \alst{h}ver standandi, &
vęifði \alst{M}jǫllni \hld\ \edtrans{\alst{m}orð-gjǫrnum}{murder-eager}{\Bfootnote{By this adjective the poet gives the Hammer something of a will of its own; cf.~\Skaldskaparmal\ 43, where the Hammer is said to always return to Thunder when thrown, perhaps by its own will.  With this may be connected the curious fact that several pagan hammer-amulets of the Wiking Age depict the Hammer with a mouth, nose, and eyes.  Three such finds are the silver hammer from Bredsättra, Öland, Sweden (item number 106666 HST), the gilt silver hammer from Ödeshög, East Geatland, Sweden (108249 HST), and the silver hammer from Baron Claes Kurck’s collection found somewhere in Scania, (106659 HST), the last of which is depicted in Fig.~\ref{fig:kurck}, p.~\pageref{fig:kurck}.}} framm, &
ok \alst{h}raun-\alst{h}vala \hld\ \alst{h}ann alla drap.\eva

\bvb He heaved off his shoulders the cauldron, standing, \\
swung the murder-eager Millner forth, \\
and all the whales of the waste \ken{ettins} he slew.\evb\evg

\begin{figure}[b]
\centering
\includegraphics[width=\textwidth]{Kurck-hammer}
\caption{Hammer-amulet with face from unknown findspot, Scania.  Wiking Age, ca. 900 \ce.  © Ola Myrin, Historiska museet/SHM. \href{https://creativecommons.org/licenses/by/4.0/deed.en}{CC BY 4.0}.  \url{https://samlingar.shm.se/media/1B654576-D251-475F-B366-6EF893A1D537}}
\label{fig:kurck}
\end{figure}

\bvg\bva\mssnote{\Regius~14v/30, \AM~6v/21}%
\Ballnote{Sts. 37--38 deal with one of Thunder’s goats going halt and an ettin farmer giving up two of his children in return.  A related story is also found in \Gylfaginning~44:9--21 and may be summarized as follows:

Thunder and Lock were on the way to visit Outyards-Lock and stayed the night with a certain farmer.  For supper Thunder cut his two goats and invited the farmer and his family to dine with him.  After they had eaten he spread the goatskins before the fire and bade the housefolk throw the bones of the goats on top of them.  The farmer’s son stealthily pried open the thigh of one of the goats and ate the marrow.  At dawn Thunder blessed the goatskins with his hammer and the goats came back to life, but the one Thelve had cut had a halt leg.  Thunder was enraged, but the farmer begged for his life and offered to give up his two children, namely his son Thelve and his daughter Wrash.  Thunder spared him, and the two became his servants.

The present stanza may reference a version of the same myth wherein Lock has a part to play in the halting of the goat, perhaps by encouraging Thelve to pry the bone open; this is strongly supported by \Skaldskaparmal\ 23 where Lock is called \emph{þjóf [\dots] hafrs} ‘thief of the he-goat’, almost certainly Thunder’s.  The farmer may be identified with Eyel, with whom the goats were left in st.~7 above, so it is probably also Eyel who gives away his two children (Thelve and Wrash) in st.~38.}%
Fóru⸗t \alst{l}ęngi, \hld\ áðr \alst{l}iggja nam &
\alst{h}afr \alst{H}lór·riða \hld\ \alst{h}alf-dauðr fyrir, &
vas \edtext{\alst{sk}ę́r}{\Afootnote{emend. from meaningless \emph{†skirr†} \Regius\AM}} \alst{sk}ǫkuls \hld\ \alst{sk}akkr ȧ bęini, &
en því hinn \alst{l}ę́-vísi \hld\ \alst{L}oki of olli.\eva

\bvb They did not journey for long before Loride’s he-goat \\
took to lie half-dead ahead. \\
The colt of the cart-pole \ken{goat} was halt in the leg, \\
but that the guile-wise Lock had caused.\evb\evg


\bvg\bva\mssnote{\Regius~14v/32, \AM~6v/22}%
En \edtrans{ér}{ye}{\Bfootnote{The listeners.  A direct address to the audience of this type is otherwise unparallelled in Eddic mythological poetry.  Such are, however, typical for the Scaldic poems with which \textlink{Hymiskvida} shares several traits; see Introduction above.  Cf.~st.~39/3--4 n.}} \alst{h}ęyrt \alst{h}afið, \hld\ \edtext{\alst{h}vęrr kann umb þat &
\alst{g}oð-mǫ́lugra}{\lemma{hvęrr \dots\ goð-mǫ́lugra ‘each god-speaking man’}\Bfootnote{Literally “each of the god-speaking ones”.  \emph{goð-mǫ́lugr} ‘god-speaking’ is an hapax, but easily understood as “learned in the lore of the gods”.}} \hld\ \alst{g}ørr at skilja, &
\alst{h}vęr af \alst{h}raun-búa \hld\ \alst{h}ann laun of fekk, &
es \alst{b}ę́ði galt \hld\ \alst{b}ǫrn sïn fyrir.\eva

\bvb But ye have heard---about that can each \\
god-speaking man more clearly discern--- \\
what recompense \emph{he} \ken*{= Thunder} got from the waste-dweller \ken{ettin} \\
when he \ken*{= the farmer} paid up both his children for it.\evb\evg


\bvg\bva\mssnote{\Regius~15r/1, \AM~6v/24}%
\edtrans{\alst{Þ}rótt-ǫflugr}{The one strong of vigour}{\Bfootnote{Thunder.}} kom \hld\ ȧ \alst{þ}ing goða &
ok \alst{h}afði \alst{h}ver, \hld\ þann’s \alst{H}ymir átti; &
\edtext{en \alst{v}éar hvęrjan \hld\ \alst{v}ęl skulu drikka &
\alst{ǫ}lðr at \alst{Ę́}gis \hld\ \edtrans{\alst{ęi}tt hǫr-męitið}{one \dots\ flax-cutter}{\Bfootnote{The latter word is an \emph{hapax} and very obscure.  It clearly belongs with \emph{hvęrjan} ‘every, each’ (m.~acc.~sg.); its nom.~sg.~form should be \emph{*hǫr-męitiðr}.  I take this as a compound formed from \emph{hǫrr} ‘flax, cord’ + \emph{męita} ‘to cut’ with an agent noun suffix \emph{-ðr}.  Since the determiner \emph{hvęrjan} implies seasonality the word is easily read as a kenning for \textsc{summer}, or specifically the time when flax is cut around late July.
This is however unprecedented amongst kennings for \textsc{summer}, which otherwise always belong to the type “mercy of the \textsc{serpent}” \parencite[109]{Meissner1921}. ---
Following the ms. evidence the pres. ed. reads \emph{ęitt} ‘one’ n. acc. sg. as modifying \emph{ǫlðr} ‘ale-feast’, but based on typical \textsc{winter}-kennings of the type “harm of the \textsc{serpent}” \textcite[46]{LaFargeGlossary} suggest the emendations \emph{ęit\emph{r}-hǫr-męiti\emph{r}} ‘venom-rope-cutter \ken{snake > winter}’ and \emph{ęit\emph{r}-*orm-\emph{m}ęi\emph{ð}i\emph{r}} ‘venom-worm-schather \ken{winter}’.}}.}%
{\lemma{en véar hvęrjan \hld\ vęl skulu drikka / ǫlðr at Ęgis \hld\ ęitt hǫr-męitið ‘but the \inx[P]{Wighers} \name{Gods} well shall drink / one ale-feast at Eagre’s every flax-cutter \ken{summer?}.’}\Bfootnote{Regardless of whether \emph{(ęitt) hǫr-męitið} means \ken{summer} or \ken{winter} (see note), these concluding lines seems to describe the Gods celebrating a yearly seasonal festival (\emph{ǫlðr}).  The fact that this is mentioned at the very end of the poem seems to give an etiological purpose to the whole narrative; with the claiming of the cauldron the Gods are henceforth able to enjoy an abundance of drink at a specified point each year.  The festival alluded to is probably equivalent to an annual agricultural festival celebrated by humans, since the cult of the Gods serves as the model for the human religion (\textlink{Voluspa}[7]/2 n.), and with that in mind we can deduce that \textlink{Hymiskvida} was likely composed for recitation either at a festival taking place in late July or at Yule---a unique insight into the performance of Eddic poetry.  This is further supported by the address of the audience in st.~38 above. ---
The word order is unusually complex for an Eddic poem but is not unprecedented in \textlink{Hymiskvida} (e.g. st.~16/3--4).  In prose word order: \emph{en véar skulu drikka vęl ęitt ǫlðr at Ę́gis hvęrjan hǫr-męitið}; or with emendation from \citeauthor{LaFargeGlossary}: \emph{en véar skulu drikka vęl ǫlðr at Ę́gis hvęrjan ęitr-hǫr-męitið.}}}\eva

\bvb The one strong of vigour came to the \inx[C]{Thing} of the Gods \\
and had the cauldron which Hymer had owned--- \\
but the \inx[P]{Wighers} \name{Gods} shall drink well \\
one ale-feast at Eagre’s every flax-cutter \ken{summer?}.\evb\evg

\sectionline
